We continue the study of linear-time transposition of $d\times d$ matrices with entries of
$b(d)=\max(1,\lceil\log_2d\rceil)$ bits on multitape Turing machines, where $L=d^2b(d)$ is the input length. Assuming an affirmative answer to the coding question of
\emph{Directed Recovery Structures in Tape Computations} (the crux), we locate where a linear-time transposer could still
exist. For oblivious machines, depth of the schedule no longer matters: a rank decomposition (hop charging) excludes every
transposer whose code merges deep data lineages only boundedly often along each path, on every schedule and with any
capacity, so the remaining oblivious question is depth reduction of the code's merge structure. For adaptive machines,
we prove a form of an earlier low-adaptivity bound that is robust against padding of bounded entry: the coarse skeleton
of the trace must carry about $L/2$ bits. Hence no linear-time transposer has bounded entry on every tape, and one whose
work tapes have bounded entry must revisit input cells unboundedly often and carry about half of the input in the
coarse motion of its input head, a condition that such revisits can meet by padding. Unconditionally, sufficiently low adaptivity reduces to the
oblivious source-cover exclusions, and every linear-time transposer, adaptive or not, yields an external-memory
transposition with linearly many I/Os that is oblivious on $2^{L-O(d)}$ inputs; an external-memory lower bound without
the indivisibility assumption would therefore refute linear-time transposition. We also prove a deep-nesting conjecture
of the series for block labels at one shift, conditional on a conjecture that the straddle counts of a random shift concentrate, and rule out one
natural building block for constructions. All proofs are written out; this is a draft, not yet refereed or formalized
in Lean.
